\begin{theorem}[Lee--Yang 边缘对数导数的 Dirichlet--resolvent 极限]\label{thm:conclusion-leyang-edge-log-derivative-dirichlet-resolvent-limit}
\leanverified{paper\_conclusion\_leyang\_edge\_log\_derivative\_dirichlet\_resolvent\_limit}
沿用结论 \ref{con:xi-terminal-zm-double-root-diffusive-sinc-limit} 的多项式族 \(Z_m(y)\)，并定义边缘缩放对数导数
\[
R_m(u):=
-\frac{1}{m^2}\,\partial_y\log Z_m(y)\Big|_{y=-u/m^2}.
\]
则对任意紧集
\[
K\Subset \CC\setminus \{4\pi^2j^2:\ j\ge 1\}
\]
成立局部一致极限
\[
R_m(u)\longrightarrow \frac{S'(u)}{S(u)}
=-\sum_{j\ge 1}\frac{1}{4\pi^2j^2-u},
\qquad u\in K,
\]
其中
\[
S(u)=\operatorname{sinc}\!\left(\frac{\sqrt u}{2}\right)
=\prod_{j\ge 1}\left(1-\frac{u}{4\pi^2j^2}\right).
\]
\end{theorem}
\begin{proof}
由结论 \ref{con:xi-terminal-zm-double-root-diffusive-sinc-limit}，
\[
\frac{Z_m(-u/m^2)}{N_m}\longrightarrow S(u)
\]
在任意紧集上一致成立，而 \(K\) 避开 \(S\) 的零点集合 \(\{4\pi^2j^2\}_{j\ge1}\)。
因此对充分大的 \(m\)，\(Z_m(-u/m^2)\) 在 \(K\) 的某个邻域内无零点，故可取解析对数并应用局部一致收敛的解析函数对数导数稳定性，得到
\[
\partial_u\log Z_m(-u/m^2)\longrightarrow \partial_u\log S(u)
\]
于 \(K\) 上局部一致。
又
\[
\partial_u\log Z_m\!\left(-\frac{u}{m^2}\right)
=
-\frac{1}{m^2}\,\partial_y\log Z_m(y)\Big|_{y=-u/m^2}
=R_m(u),
\]
故 \(R_m(u)\to S'(u)/S(u)\)。
最后，结论 \ref{con:xi-terminal-zm-diffusive-edge-weyl-spectral-determinant} 已给出 Weierstrass 乘积
\[
S(u)=\prod_{j\ge 1}\left(1-\frac{u}{4\pi^2j^2}\right),
\]
逐项取对数导数即得
\[
\frac{S'(u)}{S(u)}=-\sum_{j\ge 1}\frac{1}{4\pi^2j^2-u}.
\]
证毕。
\end{proof}

\begin{theorem}[Lee--Yang 边缘零点测度的 Dirichlet 时钟收敛]\label{thm:conclusion-leyang-edge-zero-measure-dirichlet-clock-convergence}
沿用结论 \ref{con:xi-terminal-zm-double-root-diffusive-sinc-limit} 的记号，记
\[
\rho_{m,1}>\rho_{m,2}>\cdots
\]
为 \(Z_m\) 在 \((-\infty,0)\) 上由右向左排列的零点，并定义缩放原子测度
\[
\mu_m:=\sum_{j\ge 1}\delta_{m^2(-\rho_{m,j})}.
\]
则在 \((0,\infty)\) 上有 vague 收敛
\[
\mu_m \xrightarrow{\mathrm{vague}} \mu_D:=\sum_{j\ge 1}\delta_{4\pi^2j^2}.
\]
等价地，对任意 \(\varphi\in C_c((0,\infty))\)，有
\[
\lim_{m\to\infty}\sum_{j\ge1}\varphi\!\bigl(m^2(-\rho_{m,j})\bigr)
=
\sum_{j\ge1}\varphi(4\pi^2j^2).
\]
\end{theorem}
\begin{proof}
由结论 \ref{con:xi-terminal-zm-double-root-diffusive-sinc-limit}，对每个固定 \(j\ge1\) 都有
\[
m^2(-\rho_{m,j})\longrightarrow 4\pi^2j^2.
\]
另一方面，结论 \ref{con:xi-terminal-zm-diffusive-edge-weyl-spectral-determinant} 已给出局部计数律：对任意固定 \(U>0\)，
\[
\#\{j:\ m^2(-\rho_{m,j})\le U\}
=
\left\lfloor \sqrt{\frac{U}{4\pi^2}}\right\rfloor+o(1).
\]
因此任意紧支撑测试函数 \(\varphi\) 只会看到有限多个原子，且不会发生质量逃逸。
逐个原子的收敛与局部总质量控制合并，即得所述 vague 收敛。证毕。
\end{proof}

\begin{theorem}[Lee--Yang 四次族的解析边缘与算术纤维双普适分裂]\label{thm:conclusion-leyang-analytic-edge-arithmetic-fiber-dual-universality-split}
同一四次族 \(\Pi(\lambda,y)=0\) 同时呈现两种彼此不可约化的普适极限对象：
\begin{enumerate}
\item 在边缘缩放 \(y=-u/m^2\) 下，出现由
\[
\mu_D=\sum_{j\ge1}\delta_{4\pi^2j^2}
\]
支撑的确定性 Dirichlet 时钟，其一点评价由 meromorphic resolvent
\[
u\longmapsto -\sum_{j\ge1}\frac{1}{4\pi^2j^2-u}
\]
完全编码；
\item 在良素未分歧参数上，纤维根数 \(N_p(y_0)\) 的低阶统计呈现 quartic Poisson 伪装：其概率生成函数满足
\[
\frac{1}{|U_p^{\mathrm{aff}}|}\sum_{y_0\in U_p^{\mathrm{aff}}} z^{N_p(y_0)}
=
\frac38+\frac13 z+\frac14 z^2+\frac1{24}z^4+O_z(p^{-1/2}),
\]
从而一切四次以下多项式观测都与 \(\mathrm{Pois}(1)\) 一致到 \(O(p^{-1/2})\)，而五阶下降阶乘是首个精确可见破缺阶。
\end{enumerate}
因此，该 Lee--Yang 四次族在解析边缘侧给出“无限极点、即时显影”的确定性时钟，而在算术纤维侧给出“前四阶全隐、第五阶才破”的随机 Poisson 伪装；二者共享同一代数方程，却属于两种不同的可识别机制。
\end{theorem}
\begin{proof}
第一部分正是定理 \ref{thm:conclusion-leyang-edge-log-derivative-dirichlet-resolvent-limit} 与定理 \ref{thm:conclusion-leyang-edge-zero-measure-dirichlet-clock-convergence} 的内容。
第二部分由定理 \ref{thm:conclusion-leyang-unramified-fiber-count-quartic-poisson-cloak} 与推论 \ref{cor:conclusion-leyang-fifth-falling-factorial-first-visible-break} 给出。
两侧的极限对象分别由 meromorphic resolvent 与有限阶多项式统计承载，且第二侧直到五阶才出现首个不可消去破缺，故它们不能压缩为同一种有限阶 one-point 观测机制。证毕。
\end{proof}

\begin{conjecture}[tame \texorpdfstring{$S_d$}{Sd} 覆盖的 degree-\texorpdfstring{$d$}{d} Poisson 伪装原理]\label{conj:conclusion-leyang-degree-d-poisson-cloak-principle}
设
\[
\pi:C\to \PP^1
\]
是一个 tame good reduction 的 \(S_d\)-generic 有理覆盖。
对良素 \(p\) 与未分歧参数 \(y_0\) 定义纤维根数
\[
N_p(y_0):=\#\pi^{-1}(y_0)(\FF_p).
\]
则其 good-prime 极限分布应满足：
\[
\EE\bigl[(N_p)_k\bigr]=1+O(p^{-1/2})
\qquad (1\le k\le d),
\]
而
\[
(N_p)_{d+1}\equiv 0
\]
逐点成立。
等价地，其概率生成函数在 \(z=1+t\) 处应满足
\[
\EE\bigl[(1+t)^{N_p}\bigr]
=
\sum_{k=0}^{d}\frac{t^k}{k!}+O(p^{-1/2}).
\]
换言之，good-prime 纤维根数在低阶统计上普遍伪装成参数 \(1\) 的 Poisson 律，而这种伪装的长度恰由覆盖次数 \(d\) 控制；在本节所处理的 Lee--Yang 四次族中，这一原理的 \(d=4\) 情形已被完全验证。
\end{conjecture}

\endinput
